\color{unidarkgreen}{\subsection{ Устройство резервирования отказа выключателя}}
\color{black}

\begin{enumerate}

\item
Устройство резервирования отказа выключателя (УРОВ) предназначено для выполнения функции ближнего резервирования защит. В случае отказа выключателя повреждённого присоединения УРОВ обеспечивает отключение смежных выключателей с целью селективного отделения повреждённого участка электрической сети.

\item
В случае присоединения линии к системе через два выключателя для каждого из них можно задействовать свою схему УРОВ. Далее приведено описание УРОВ для выключателя В1, логика УРОВ для второго выключателя аналогична.

\item
В устройстве применяется индивидуальный принцип реализации функции УРОВ. Пуск УРОВ происходит при срабатывании внутренних защит, выявлении неполнофазного режима (опционально, разрешается при введенном программном ключе \textbf{«Пуск УРОВ В1 от ЗНР»}), либо от внешних защит. 

\item
Для выявления факта отказа выключателя предусмотрен контроль наличия тока через выключатель. Контроль осуществляется измерительным органом УРОВ, порог срабатывания которого задается уставкой \textbf{«Iср УРОВ В1»}.

\item
Для обеспечения надежного пуска УРОВ при кратковременном пропадании сигналов пуска до полного обесточивания отказавшего выключателя предусмотрен подхват от ИО УРОВ, ввод которого осуществляется с помощью программного ключа \textbf{«Фикс. УРОВ В1»}, сброс пуска произойдет только при возврате ИО УРОВ.

\item
Входной сигнал ФСУ \textit{«УРОВ В1/УРОВ: Пуск УРОВ НН»} предусмотрен для пуска УРОВ от защит, срабатывание которых не сопровождается увеличением тока или сопровождается протеканием малых токов КЗ, недостаточных для срабатывания токовых ИО УРОВ. Например, от пускового органа УРОВ НН, предусматриваемого в защитах (авто)трансформатора, который срабатывает при КЗ на стороне НН (авто)трансформатора за токоограничивающим реактором. Данный пуск осуществляется с контролем включенного положения выключателя.

\item
Возможны два режима работы УРОВ:
\begin{itemize}
    \item с автоматической проверкой исправности выключателя;
    \item с дублированным пуском.
\end{itemize}

\item
В режиме «с автоматической проверкой исправности выключателя» программный ключ \textbf{«Режим УРОВ В1 с дубл. пуском»} устанавливается в положение \textit{«Выведено»}, а программный ключ \textbf{«УРОВ В1 на себя»} -- \textit{«Введено»}. Данный режим выполняет функцию резервирования при обрыве цепи действия на отключение выключателя от защит. При появлении пуска УРОВ выдается команда на отключение «своего» выключателя, что позволяет исключить ложное срабатывание УРОВ на смежные выключатели. Выдержка времени действия 
УРОВ на «свой» выключатель регулируется уставкой \textbf{«Тср УРОВ В1 на себя»}. Если же отключения выключателя не происходит, значит неисправна не цепь от защит к ЭМО, а сам ЭМО. Тогда с проверкой срабатывания ИО УРОВ и с большей выдержкой времени \textbf{«Тср УРОВ В1»} формируется сигнал срабатывания УРОВ на смежные выключатели. Контроль наличия тока при действии УРОВ «на себя» может быть выведено программным ключом \textbf{«Контр. по I УРОВ В1 на себя»}.

\item
Действие УРОВ на себя осуществляется:
\begin{itemize}
    \item от сигналов внешнего пуска \textit{«УРОВ В1/УРОВ: Пуск УРОВ №1»} и \textit{«УРОВ В1/УРОВ: Пуск УРОВ №2»}, для которых можно ввести контроль по току и подхват от ИО УРОВ; 
    \item от сигнала \textit{«УРОВ В1/УРОВ: Пуск УРОВ НН»} с контролем положения выключателя, без контроля по току и без возможности подхвата от ИО УРОВ.
\end{itemize}

\item
В режиме «с дублированным пуском» программный ключ \textbf{«Режим УРОВ В1 с дубл. пуском»} устанавливается в положение \textit{«Введено»}, а программный ключ \textbf{«УРОВ В1 на себя»} – \textit{«Выведено»}. В данном режиме пуск УРОВ разрешается только при срабатывании защит с действием на отключение – пропадание входного сигнала ФСУ \textbf{«Контроль В: Готовность ЦО»} является признаком срабатывания защит и работе цепи отключения. С выдержкой времени \textbf{«Тср УРОВ В1»} формируется сигнал срабатывания УРОВ на смежные выключатели.

\item
Действие УРОВ на смежные выключатели осуществляется:
\begin{itemize}
    \item от сигналов внешнего пуска \textit{«УРОВ В1/УРОВ: Пуск УРОВ №1»} и \textit{«УРОВ В1/УРОВ: Пуск УРОВ №2»}, от срабатывания ЗНР, от срабатывания внутренних защит. Для всех перечисленных сигналов можно ввести подхват от ИО УРОВ; 
    \item от сигнала \textit{«УРОВ В1/УРОВ: Пуск УРОВ НН»} с контролем положения выключателя, без контроля по току и без возможности подхвата от ИО УРОВ.
\end{itemize}

\item
УРОВ В1 может быть оперативно выведено из работы от внешнего сигнала \textit{«УРОВ В1/УРОВ: Вывод»}, назначенного на соответствующий дискретный вход или функциональную клавишу.

\item 
Функциональная схема логики  УРОВ В1 приведена на рисунке \ref{pict_urov:UROV}. Функциональная схема логики УРОВ В2 выполняется аналогично схеме логики УРОВ В1.

\medskip

\begin{figure}[H]
\centering
\includegraphics[width=1\textwidth,height=1\textheight,keepaspectratio]{\fbpath/03.01 ЛВ/761.3311 УРОВ/img/UROV.pdf}
\captionof{figure}{Функциональная схема логики УРОВ В1}\label{pict_urov:UROV}
\end{figure}

\medskip

\item
\textls[-10]{Уставки УРОВ В1 представлены в таблице \ref{app_set:urov_q1}. Уставки УРОВ В2 представлены в таблице \ref{app_set:urov_q2}.}

\end{enumerate}